% Standalone problem document, generated from the adjacent README.md.
% Compile with XeLaTeX. Mathematical content is maintained in Markdown.
\documentclass[11pt,a4paper]{article}
\usepackage[fontsize=10.5pt]{scrextend}
\usepackage[margin=22mm,headheight=15pt,headsep=9mm,footskip=11mm]{geometry}
\usepackage{fontspec}
\setmainfont{texgyrepagella}[Extension=.otf,UprightFont=*-regular,BoldFont=*-bold,ItalicFont=*-italic,BoldItalicFont=*-bolditalic]
\setsansfont{texgyreheros}[Extension=.otf,UprightFont=*-regular,BoldFont=*-bold,ItalicFont=*-italic,BoldItalicFont=*-bolditalic]
\setmonofont{texgyrecursor}[Extension=.otf,UprightFont=*-regular,BoldFont=*-bold,ItalicFont=*-italic,BoldItalicFont=*-bolditalic,Scale=MatchLowercase]
\usepackage{amsmath,amssymb}
\usepackage{unicode-math}
\setmathfont{texgyrepagella-math.otf}
\usepackage{xcolor,microtype,enumitem,needspace,fancyhdr}
\usepackage{bookmark}
\definecolor{ink}{HTML}{17344B}
\definecolor{accent}{HTML}{197D80}
\definecolor{muted}{HTML}{596773}
\hypersetup{colorlinks=true,linkcolor=accent,urlcolor=accent,pdftitle={AC-10 - An
explicit rational Valiant-rigid matrix
family},pdfauthor={Open Problems in Numerical Linear Algebra}}
\urlstyle{same}
\setlength{\parindent}{0pt}
\setlength{\parskip}{5pt plus 1pt minus 1pt}
\linespread{1.04}
\setlength{\emergencystretch}{3em}
\widowpenalty=10000
\clubpenalty=10000
\displaywidowpenalty=10000
\allowdisplaybreaks[1]
\setlist{nosep,leftmargin=1.5em}
\providecommand{\tightlist}{\setlength{\itemsep}{0pt}\setlength{\parskip}{0pt}}
\providecommand{\pandocbounded}[1]{#1}
\pagestyle{fancy}
\fancyhf{}
\fancyhead[L]{\sffamily\footnotesize\color{muted}OPEN PROBLEMS IN NUMERICAL LINEAR ALGEBRA}
\fancyhead[R]{\sffamily\bfseries\footnotesize\color{accent}AC-10}
\fancyfoot[L]{\parbox[b]{0.9\textwidth}{\sffamily\fontsize{8}{10}\selectfont\color{muted}Editorial ratings; a bounded search cannot certify that a problem remains unsolved.\\Literature check: 2026-09-08}}
\fancyfoot[R]{\sffamily\footnotesize\color{muted}\thepage}
\renewcommand{\headrulewidth}{0.3pt}
\renewcommand{\headrule}{\hbox to\headwidth{\color{accent}\leaders\hrule height \headrulewidth\hfill}}
\makeatletter
\renewcommand{\section}{\Needspace{5\baselineskip}\@startsection{section}{1}{0pt}{12pt plus 2pt minus 2pt}{4pt}{\sffamily\bfseries\large\color{ink}}}
\renewcommand{\subsection}{\Needspace{4\baselineskip}\@startsection{subsection}{2}{0pt}{9pt plus 1pt minus 1pt}{3pt}{\sffamily\bfseries\normalsize\color{ink}}}
\makeatother
\setcounter{secnumdepth}{0}
\begin{document}
{\sffamily\small\color{muted}Arithmetic and complexity\par}
\vspace{3pt}
{\raggedright\hyphenpenalty=10000\exhyphenpenalty=10000\sffamily\bfseries\fontsize{21}{24}\selectfont\color{ink}An
explicit rational Valiant-rigid matrix family\par}
\vspace{7pt}
{\sffamily\small\textbf{Difficulty:} extreme\quad\textbf{Importance:} broadly
interesting\par}
\vspace{4pt}
{\color{accent}\hrule height 0.8pt}
\vspace{6pt}
\textbf{Provenance:} rational-field specialization of a source-stated
explicit-rigidity target.

\textbf{Status:} no construction meeting this target found in screening
on 2026-09-08.

For \(A\in\mathbb Q^{n\times n}\) and an integer \(0\leq r\leq n\),
define its rigidity over \(\mathbb Q\) by \[
R_A^{\mathbb Q}(r)
=\min\left\{
|\{(i,j):A_{ij}\ne B_{ij}\}|:
B\in\mathbb Q^{n\times n},\
\operatorname{rank}_{\mathbb Q}(B)\leq r
\right\}.
\] Construct a constant \(\delta>0\) and a family
\(A_n\in\mathbb Q^{n\times n}\) such that a single deterministic
algorithm, on input the unary encoding \(1^n\), outputs all rational
entries of \(A_n\) in binary using polynomially many bit operations in
\(n\), and \[
R_{A_n}^{\mathbb Q}
\left(\left\lfloor\frac{n}{\log\log n}\right\rfloor\right)
\geq n^{1+\delta}
\] for every sufficiently large integer \(n\). Here \(\log\) denotes the
natural logarithm. The output convention requires polynomial total bit
length; an exponentially long algebraic description does not meet this
requirement.

Rigidity measures how resistant a matrix is to rank reduction by sparse
entry changes. An explicit family at this scale would supply hard
instances for linear circuit computation. The source states the target
over a field; this entry explicitly selects rational entries and
rational perturbations, rather than asserting that all choices of field
are equivalent.

\subsection{References}\label{references}

\begin{enumerate}
\def\labelenumi{\arabic{enumi}.}
\tightlist
\item
  J. Alman and J. Liang, \emph{Low Rank Matrix Rigidity: Tight Lower
  Bounds and Hardness Amplification}, arXiv:2502.19580 (2025), Section 1
  and Section 1.1: definition of rigidity, uniform polynomial-time
  explicitness, and the \(n/\log\log n\) Valiant target.
  \href{https://arxiv.org/html/2502.19580v1}{Paper}.
\item
  L. Hambardzumyan, K. Myasnikov, A. Riazanov, M. Shirley, and A.
  Shraibman, \emph{Spiky Rank and Its Applications to Rigidity and
  Circuits}, ICALP 2026, LIPIcs 374, article 106, Section 1.2.1:
  explicit rigidity remains insufficient for the relevant circuit lower
  bounds.
  \href{https://drops.dagstuhl.de/entities/document/10.4230/LIPIcs.ICALP.2026.106}{Published
  paper}.
\item
  L. G. Valiant, \emph{Graph-theoretic arguments in low-level
  complexity}, MFCS 1977, pp.~162--176, the original rigidity connection
  to linear circuits.
  \href{https://doi.org/10.1007/3-540-08353-7_135}{Article}.
\item
  Z. Dvir and A. Liu, \emph{Fourier and Circulant Matrices are Not
  Rigid}, arXiv:1902.07334 (2019), main non-rigidity results for
  prominent structured candidates.
  \href{https://arxiv.org/abs/1902.07334}{Paper}.
\end{enumerate}

\subsection{Status check --- 2026-09-08}\label{status-check-2026-09-08}

checked 2502.19580v1 (2025-02-26), the July 2026 published paper, and
searches ``matrix rigidity explicit rational 2026 Valiant construction''
and ``Valiant rigidity solved''. The 2025 bounds for logarithmic target
rank and its conditional amplification results do not give the displayed
higher-rank target. Known non-rigidity results prevent treating Walsh or
Fourier matrices as established candidates. No qualifying
polynomial-time rational construction or resolution was found. This
status applies to the stated rational specialization.
\end{document}
